% SPDX-License-Identifier: CC-BY-SA-4.0
% Author: Matthieu Perrin
% Part: <Nom de la partie>
% Section: <Nom de la section>
% Sub-section: <Nom de la sous-section>  % (facultatif, laisser vide si non utilisé)
% Frame: <Titre de la slide>

\begingroup

\SetKwFunction{IsPow}{isPow2}
\SetKwFunction{AlgoSearchPerfect}{search\_perfect}

\begin{frame}{Généralisation de la preuve d'indécidabilité}

  \onBlock[top=-5mm]{Définition -- Réduction par mappage}{
    Soient $A$ et $B$ deux problèmes de décision.
    
    Une \structure{réduction de $A$ vers $B$} est une fonction
    \alert{$f : \mathcal{I}(A) \rightarrow \mathcal{I}(B)$} telle que
    \begin{itemize}
    \item $f$ est \alert{calculable} ;
    \item $\forall u \in \mathcal{I}(A),\quad \alert{u\in \textsc{pos}(A) \Leftrightarrow f(u)\in \textsc{pos}(B)}$.
    \end{itemize}
  }

  \obExampleBlock<1>[anchor=north, y=12mm]{Exemple -- De \textsc{Universal} vers \textsc{Entscheidungsproblem}}{

    \vspace{2mm}
    
    $$
    f: \left\{\begin{array}{rcl}
    \mathcal{I}(\textsc{Universal}) & \rightarrow & \mathcal{I}(\textsc{Entscheidungsproblem})\\
    \langle P, u \rangle & \mapsto & ``u\in \llbracket P \rrbracket"
    \end{array}\right.
    $$

    \vspace{3mm}

    \begin{itemize}
    \item $f$ est calculable :
      \begin{itemize}
      \item simple encodage syntaxique
      \end{itemize}
    \item $f$ préserve la positivité :
      \begin{itemize}
      \item par définition de \textsc{Universal}
      \end{itemize}
    \end{itemize}
  }
  
  \onBlock<2->[anchor=north, y=12mm]{Intuition -- $B$ est au moins aussi difficile que $A$}{}
  
  \on<2->[y=-16mm]{
    \begin{tikzpicture}
      \begin{scope}
        \clip (0,0) ellipse (2cm and 1cm);

        \begin{scope}
          \clip (-2,0) rectangle (2,1);
          \fill[structure!20] (-2,0) rectangle (2,1);
          \pgfdeclareradialshading{alertToStructure}{\pgfpoint{0cm}{0cm}}{color(0cm)=(alert!20);color(6mm)=(alert!20);color(9mm)=(structure!20);color(1cm)=(structure!20)}
          \shade[shading=alertToStructure] (2,0) ellipse (2.5cm and .75cm);
        \end{scope}
        \begin{scope}
          \clip (-2,0) rectangle (2,-1);
          \fill[example!20] (-2,0) rectangle (2,-1);
          \pgfdeclareradialshading{alertToExample}{\pgfpoint{0cm}{0cm}}{color(0cm)=(alert!20);color(6mm)=(alert!20);color(9mm)=(example!20);color(1cm)=(example!20)}
          \shade[shading=alertToExample] (2,0) ellipse (2.5cm and .75cm);
        \end{scope}
        
        \draw[black!40,thick] (0,0) ellipse (2cm and 1cm);
        \draw[black!40] (-2,0) -- (2,0);
        \scriptsize
        \node[align=center, structure] at (0, .8) {$\textsc{pos}(A)$};
        \node[align=center, example]   at (0,-.8) {$\textsc{neg}(A)$};
      \end{scope}
      
      \begin{scope}
        \clip (5,0) ellipse (2cm and 1cm);
        \begin{scope}
          \clip (3,0) rectangle (7,1);
          \fill[structure!20] (3,0) rectangle (7,1);
          \pgfdeclareradialshading{alertToStructure}{\pgfpoint{0cm}{0cm}}{color(0cm)=(alert!20);color(6mm)=(alert!20);color(9mm)=(structure!20);color(1cm)=(structure!20)}
          \shade[shading=alertToStructure] (7,0) ellipse (2.5cm and .75cm);
        \end{scope}
        \begin{scope}
          \clip (3,0) rectangle (7,-1);
          \fill[example!20] (3,0) rectangle (7,-1);
          \pgfdeclareradialshading{alertToExample}{\pgfpoint{0cm}{0cm}}{color(0cm)=(alert!20);color(6mm)=(alert!20);color(9mm)=(example!20);color(1cm)=(example!20)}
          \shade[shading=alertToExample] (7,0) ellipse (2.5cm and .75cm);
        \end{scope}
        \draw[black!40,thick] (5,0) ellipse (2cm and 1cm);
        \draw[black!40] (3,0) -- (7,0);
        \scriptsize
        \node[align=center, structure] at (4, .6) {$\textsc{pos}(B)$};
        \node[align=center, example]   at (4,-.6) {$\textsc{neg}(B)$};
      \end{scope}

      \begin{scope}
        \clip (4.5,0) ellipse (1cm and .5cm);
        \begin{scope}
          \clip (3,0) rectangle (7,1);
          \fill[structure!40] (3,0) rectangle (7,1);
          \pgfdeclareradialshading{alertToStructure}{\pgfpoint{0cm}{0cm}}{color(0cm)=(alert!40);color(6mm)=(alert!40);color(9mm)=(structure!40);color(1cm)=(structure!40)}
          \shade[shading=alertToStructure] (7,0) ellipse (2.5cm and .75cm);
        \end{scope}
        \begin{scope}
          \clip (3,0) rectangle (7,-1);
          \fill[example!40] (3,0) rectangle (7,-1);
          \pgfdeclareradialshading{alertToExample}{\pgfpoint{0cm}{0cm}}{color(0cm)=(alert!40);color(6mm)=(alert!40);color(9mm)=(example!40);color(1cm)=(example!40)}
          \shade[shading=alertToExample] (7,0) ellipse (2.5cm and .75cm);
        \end{scope}
        \draw[black!40,thick] (4.5,0) ellipse (1cm and .5cm);;
        \draw[black!40] (3,0) -- (7,0);
        \scriptsize
        \node[align=center, structure] at (4.5, .2) {$f(\textsc{pos}(A))$};
        \node[align=center, example]   at (4.5,-.2) {$f(\textsc{neg}(A))$};
      \end{scope}

      \footnotesize
      \node[inner sep=0pt, alert]     (Ah1) at (1.85,0)    {$\bullet$};
      \node[inner sep=0pt, structure] (Ap1) at (-1.5, .5)  {$\bullet$};
      \node[inner sep=0pt, example]   (An1) at (-1.5,-.5)  {$\bullet$};
      \node[inner sep=0pt, structure] (Ap2) at (1.1, .7)   {$\bullet$};
      \node[inner sep=0pt, example]   (An2) at (1.1,-.7)   {$\bullet$};
      
      \node[inner sep=0pt, alert]     (Bh1) at (5.4,0)     {$\bullet$};
      \node[inner sep=0pt, structure] (Bp1) at (3.75, .25) {$\bullet$};
      \node[inner sep=0pt, example]   (Bn1) at (3.75,-.25) {$\bullet$};
      \node[inner sep=0pt, structure] (Bp2) at (5, .35)    {$\bullet$};
      \node[inner sep=0pt, example]   (Bn2) at (5,-.35)    {$\bullet$};

      \node[inner sep=0pt, alert]     (Bh3) at (5.6,0)     {$\bullet$};
      \node[inner sep=0pt, structure] (Bp3) at (5.15, .5)  {$\bullet$};
      \node[inner sep=0pt, example]   (Bn3) at (5.15,-.5)  {$\bullet$};
      
      \tiny
      \path[-latex, alert]     (Ah1.center) edge             (Bh1) ; 
      \path[-latex, structure] (Ap1.center) edge[bend left]  (Bp1) ; 
      \path[-latex, example]   (An1.center) edge[bend right] (Bn1) ; 
      \path[-latex, structure] (Ap2.center) edge[bend left]  (Bp2) ; 
      \path[-latex, example]   (An2.center) edge[bend right] (Bn2) ; 
      
      \node[alert, below left]      at (1.85,0)           {$\langle \AlgoSearchPerfect, 3 \rangle$};
      \node[structure, below right] at (-1.5, .5)         {$\langle \IsPow, 8 \rangle$};
      \node[example, above right]   at (-1.5,-.5)         {$\langle \IsPow,12 \rangle$};
      \node[structure, below left]  at (1.1, .7)          {$\langle \AlgoSearchPerfect, 2 \rangle$};
      \node[example, above left]    at (1.1,-.7)          {$\langle \IsPow, 0 \rangle$};
      \node[below right, alert]     at ([xshift=-1mm]Bh3) {$\exists n$ parfait impair};
      \node[above, structure]       at ([xshift=2mm]Bp3)  {$1+2+3=6$};
      \node[below, example]         at (Bn3)              {$\exists k, 0 = 2^k$};

      \normalsize
      \node at (0,1.5)    {$A = \textsc{Universal}$};
      \node at (5,1.5)    {$B = \textsc{Entscheidungsproblem}$};
      \node at (2.5,-1.5) {$f: \langle P, u \rangle \mapsto ``u\in \llbracket P \rrbracket" $};

    \end{tikzpicture}
  }

  \on<2->[text, bottom=-2mm]{
    \begin{itemize}
    \item Toute instance difficile de $A$ reste difficile après réduction dans $B$
    \item $B$ contient des instances au moins aussi difficiles que $A$ 
    \end{itemize}
  }
  
\end{frame}

\endgroup
